% Documentary recovery for Article V; RESULTADO_RECUPERADO.
% Source owner: /Users/ruben/Documents/New project/output/REVISION_EDITORIAL_HMT_MD_20260905/01_FUENTES/INTEGRAL/tree/New project/output/ENSAMBLADO_CAUSAL_RESTAURADO_HMT_MD_20260905/fuente/manuscrito/integracion_83/parte_iv/body/B014_ch55_estadisticas_ocupacion_7549b83967f5_04d_ocupacion_estadistica.tex
% Source lines: 1--225.
% Only reference labels and, where appropriate, heading levels are adapted.
% The recovery matrix preserves the correspondence and dependencies.
\section{Microcanonical enumeration of states on a finite screen}
\label{v:rec:ocupacion:section:01}

Consider the finite realization
\[
 \mathcal H_{\rm occ}=\bigoplus_{k=0}^{3}\CC^{g_k},
 \qquad
 H_{\rm occ}|_{\CC^{g_k}}=E_k^{\rm occ} I,
\]
with levels and multiplicities
\[
 E^{\rm occ}=(0,3,6,9),\qquad g=(1,6,12,8),
 \qquad(N,E)=(12,66).
\]
The sum \(1+6+12+8=27\) makes \(\mathcal H_{\rm occ}\) a
twenty-seven-mode screen. This spectrum is the data of a specified finite
representation; its cardinality does not suffice to identify it with any other
occurrence of \(27\) in HMT\@. The morphism relating a particular
boundary screen to \((\mathcal H_{\rm occ},H_{\rm occ})\) must specify both
the fibers and the spectral operator.
The fermionic and bosonic generating functions are
\[
 F(y,x)=\prod_k(1+yx^{E_k^{\rm occ}})^{g_k},
 \qquad
 B(y,x)=\prod_k(1-yx^{E_k^{\rm occ}})^{-g_k}.
\]

\begin{proposition}[Microcanonical cardinalities]
\label{v:rec:prop:cardinalidades-microcanonicas}
\[
 [y^{12}x^{66}]F=2\,104\,494,
 \qquad
 [y^{12}x^{66}]B=259\,436\,430.
\]
\end{proposition}
\begin{proof}
In the fermionic product, choosing \(r\) of the \(g_k\) modes contributes
\(\binom{g_k}{r}y^rx^{rE_k^{\rm occ}}\). In the bosonic product, distributing \(r\)
indistinguishable occupations among \(g_k\) modes contributes
\(\binom{g_k+r-1}{r}y^rx^{rE_k^{\rm occ}}\). Finite convolution over the
four levels yields the stated integers.
\end{proof}

\section{Bose--Einstein and Fermi--Dirac distributions}
\begingroup\fontsize{10}{11.5}\selectfont\setlength{\parskip}{1.5pt}
\setlength{\abovedisplayskip}{3pt}\setlength{\belowdisplayskip}{3pt}
\setlength{\abovedisplayshortskip}{2pt}\setlength{\belowdisplayshortskip}{2pt}
\label{v:rec:ocupacion:section:02}

The equilibrium formulation starts from the Hamiltonian and the total
occupation number,
\begin{equation}
 H=\sum_kE_k^{\rm occ} n_k,
 \qquad
 N=\sum_k n_k,
 \qquad
 K_\mu:=H-\mu N.
 \label{v:rec:eq:hamiltoniano-gran-canonico-rev3}
\end{equation}
For \(\beta>0\), the grand canonical state is
\begin{equation}
 \rho_{\beta,\mu}
 =\frac{e^{-\beta K_\mu}}{\mathcal Z_{\beta,\mu}},
 \qquad
 \mathcal Z_{\beta,\mu}
 =\operatorname{Tr}(e^{-\beta K_\mu}).
 \label{v:rec:eq:estado-gran-canonico-rev3}
\end{equation}

\begin{theorem}[Grand canonical factorization and occupation laws]
\label{v:rec:thm:factorizacion-gran-canonica-rev3}
For energy levels \(E_k^{\rm occ}\) with degeneracies \(g_k\), the
fermionic trace factorizes as
\begin{equation}
 \mathcal Z_{\rm FD}
 =\prod_k\bigl(1+e^{-\beta(E_k^{\rm occ}-\mu)}\bigr)^{g_k},
 \qquad
 \langle n_k\rangle_{\rm FD}
 =\frac{g_k}{e^{\beta(E_k^{\rm occ}-\mu)}+1}.
 \label{v:rec:eq:factorizacion-fd-rev3}
\end{equation}
In the bosonic sector, under the condition
\(\mu<\inf_kE_k^{\rm occ}\),
\begin{equation}
 \mathcal Z_{\rm BE}
 =\prod_k\bigl(1-e^{-\beta(E_k^{\rm occ}-\mu)}\bigr)^{-g_k},
 \qquad
 \langle n_k\rangle_{\rm BE}
 =\frac{g_k}{e^{\beta(E_k^{\rm occ}-\mu)}-1}.
 \label{v:rec:eq:factorizacion-be-rev3}
\end{equation}
\end{theorem}

\begin{proof}
The operator \(K_\mu\) is a sum of commuting occupation operators, so
its trace factorizes over modes. In each fermionic mode, the
occupations \(0\) and \(1\) are summed, and the corresponding factor is
\(1+u_k\), where
\(u_k=e^{-\beta(E_k^{\rm occ}-\mu)}\). In each bosonic mode, all
nonnegative occupations are summed, giving the geometric series
\((1-u_k)^{-1}\); the condition on \(\mu\) ensures \(u_k<1\). The
powers \(g_k\) incorporate the degeneracy. Finally,
\[
 \langle n_k\rangle
 =-\frac1\beta\frac{\partial}{\partial E_k^{\rm occ}}
   \log\mathcal Z
\]
yields the two occupation laws
\cite{Bose1924,Einstein1924,Fermi1926,Dirac1926,Pathria2011}.
\end{proof}

The microcanonical derivation recovers the same result in the
thermodynamic regime. For mean fermionic occupations
\(p_k=m_k/g_k\), the combinatorial entropy per level is
\begin{equation}
 S_k\simeq
 g_k\bigl[-p_k\log p_k-(1-p_k)\log(1-p_k)\bigr].
 \label{v:rec:eq:entropia-fd-nivel-rev3}
\end{equation}
Variation of \(\sum_kS_k\), subject to
\(\sum_kg_kp_k=N\) and
\(\sum_kg_kE_k^{\rm occ}p_k=E\), gives
\begin{equation}
 \log\frac{1-p_k}{p_k}=\lambda_N+\beta E_k^{\rm occ},
 \qquad
 \mu=-\frac{\lambda_N}{\beta},
 \label{v:rec:eq:variacion-fd-rev3}
\end{equation}
and hence the Fermi--Dirac occupation in
\eqref{v:rec:eq:factorizacion-fd-rev3}. The multipliers \(\beta\) and \(\mu\)
are the dual coordinates of the energy and population constraints.
The symbol \(\lambda_N=-\beta\mu\) denotes the multiplier associated with
the occupation number; throughout this article, \(\alpha\) remains reserved
for the fine-structure constant.

For mean bosonic occupations
\(\overline n_k=m_k/g_k\), the logarithm of the level multiplicity is
\begin{equation}
 S_k\simeq
 g_k\bigl[(1+\overline n_k)\log(1+\overline n_k)
          -\overline n_k\log\overline n_k\bigr].
 \label{v:rec:eq:entropia-be-nivel-rev3}
\end{equation}
Variation of \(\sum_kS_k\), subject to
\(\sum_kg_k\overline n_k=N\) and
\(\sum_kg_kE_k^{\rm occ}\overline n_k=E\), yields
\begin{equation}
 \log\frac{1+\overline n_k}{\overline n_k}
 =\lambda_N+\beta E_k^{\rm occ},
 \qquad
 \overline n_k
 =\frac{1}{e^{\beta(E_k^{\rm occ}-\mu)}-1},
 \qquad
 \mu=-\frac{\lambda_N}{\beta}.
 \label{v:rec:eq:variacion-be-rev3}
\end{equation}
This is the Bose--Einstein occupation in
\eqref{v:rec:eq:factorizacion-be-rev3}. The condition
\(\mu<\inf_kE_k^{\rm occ}\) ensures positivity of the denominator and
convergence of each modal factor.

In the finite instance constructed here, the two constraints determine
\[
 (\beta_{\rm FD},\mu_{\rm FD})
 =(0.15307011353415928,\,4.502865197597043),
\]
\[
 (\beta_{\rm BE},\mu_{\rm BE})
 =(0.05456727960068316,\,-15.92569781906737).
\]
If \(m_k^{\rm mic}\) denotes the mean occupation obtained by exact
microcanonical enumeration and \(m_k^{\rm eq}\) the evaluation of the
corresponding distribution, the errors
\(\sum_k|m_k^{\rm mic}-m_k^{\rm eq}|\) are
\[
 0.09897380039002812\quad\hbox{and}\quad
 0.9190625715582017,
\]
respectively. The enumeration traverses all occupation quadruples,
calculates their integer multiplicities, and solves the two constraint
equations. Both laws are realized on the same screen. The physical
choice between them requires a grading and a Hamiltonian compatible with the
trajectory representation. The thermal parameter is
\(\beta=(k_B\Theta)^{-1}\), with the dimensional transduction developed in
Section~\ref{v:ant:sec:ii-kb-aut-desarrollo}.

\par\endgroup
\section{Quantified convergence to the Maxwell--Boltzmann distribution}
\label{v:rec:ocupacion:section:03}

The relation between the two quantum statistics and the classical regime can
be formulated without recourse to an informal approximation. For each mode, let
\[
 u_k:=e^{-\beta(E_k^{\rm occ}-\mu)}.
\]
In the graded-carrier specialisation of
Section~\ref{v:sec:capacidad-espectral-termica}, the excited degrees
$n=1,2$ carry $E_n=n\epsilon_0$, $\mu=0$, and
$\beta=\tau^{-1}$; therefore $u_n=x^n$, with
$x=e^{-\epsilon_0/\tau}$, and their multiplicities are $380$ and $649$.
Degree $n=0$ remains the marked origin. Thus
\eqref{v:eq:particion-portador-termico} supplies the finite trace before
completion, whereas the Bose--Einstein and Fermi--Dirac formulae remain
the occupation laws of the completed modes. In particular, the bounds
below apply to the excited sectors as $x\to0$.
The occupations per state are written
\[
 \overline n_k^{\rm FD}=\frac{u_k}{1+u_k},
 \qquad
 \overline n_k^{\rm BE}=\frac{u_k}{1-u_k};
\]
the multiplicities \(g_k\) are incorporated afterward through
\(\langle n_k\rangle=g_k\overline n_k\).

\begin{theorem}[Common classical limit of quantum statistics]
\label{v:rec:thm:limite-clasico-fd-be-rev2}
If \(0\leq u_k\leq\delta<1\), then
\[
 \left|\overline n_k^{\rm FD}-u_k\right|\leq u_k^2,
 \qquad
 \left|\overline n_k^{\rm BE}-u_k\right|
 \leq\frac{u_k^2}{1-\delta}.
\]
Consequently, in the dilute regime \(u_k\to0\),
\[
 \boxed{
 \overline n_k^{\rm FD}\sim
 \overline n_k^{\rm BE}\sim
 e^{-\beta(E_k^{\rm occ}-\mu)}} ,
\]
which is the Maxwell--Boltzmann distribution per mode.
\end{theorem}

\begin{proof}
The exact identities
\[
 u_k-\frac{u_k}{1+u_k}=\frac{u_k^2}{1+u_k},
 \qquad
 \frac{u_k}{1-u_k}-u_k=\frac{u_k^2}{1-u_k}
\]
give the bounds. For \(0<u_k\leq\delta\), dividing by \(u_k\) and
taking \(u_k\downarrow0\) yields
the asymptotic equivalence. The first term of order \(u_k^2\) retains the
physical difference: exclusion decreases fermionic occupation, whereas
accumulation increases bosonic occupation.
\end{proof}
